\chapter{Technical Questions}
\label{ch:technical}


For candidates preparing for technical roles such as \textbf{Assistant Director (ICT), Assistant Programmer, Assistant Maintenance Engineer, Assistant Database Administrator, Assistant Network Engineer, or similar IT-related positions}, preparation should extend well beyond rote memorization. Interview boards frequently transition from fundamental definitions into deeper conceptual inquiries, requiring candidates to evaluate architectural trade-offs, compare competing concepts, or solve practical problems on the spot.

Below is a \textbf{large viva bank of 300+ questions}, organized from fundamentals to advanced topics. The answers are deliberately \textbf{short, intuitive, and viva-ready} rather than textbook-length.

\medskip

\section{Assistant Programmer Viva: 300+ CSE Questions \& Answers}

\subsection{How to use this}

For each topic, learn in three layers:

\begin{enumerate}
  \item \textbf{Definition} --- one sentence.
  \item \textbf{Intuition} --- why it exists / how it works.
  \item \textbf{Follow-up} --- comparison, example, complexity, or implementation.
\end{enumerate}

In the actual viva, give the \textbf{shortest correct answer first}. Elaborate only if the board asks.

\medskip

\section{Computer Fundamentals}

\subsection{What is a computer?}

A computer is an electronic programmable machine that accepts data as input, processes it according to instructions, stores it, and produces information as output.

The basic cycle is:

\textbf{Input $\rightarrow$ Processing $\rightarrow$ Output $\rightarrow$ Storage}

\medskip

\subsection{What is the difference between data and information?}

\textbf{Data} is raw facts.

\textbf{Information} is processed data that has meaningful context.

Example:

\begin{lstlisting}
Data:       85, 72, 91
Information: Average mark = 82.67
\end{lstlisting}

\medskip

\subsection{What is hardware?}

Hardware refers to the physical components of a computer.

Examples:

\begin{itemize}
  \item CPU
  \item RAM
  \item motherboard
  \item keyboard
  \item SSD
  \item monitor
\end{itemize}

\medskip

\subsection{What is software?}

Software is a collection of programs and instructions that tell hardware what to do.

Examples:

\begin{itemize}
  \item Windows
  \item Linux
  \item Chrome
  \item MySQL
  \item Python
\end{itemize}

\medskip

\subsection{What are the major types of software?}

Two broad categories:

\begin{enumerate}
  \item \textbf{System software}
  \item \textbf{Application software}
\end{enumerate}

Operating systems, device drivers and compilers are system software.

Word processors, browsers and accounting applications are application software.

\medskip

\subsection{What is firmware?}

Firmware is software stored in non-volatile memory that provides low-level control of hardware.

Examples:

\begin{itemize}
  \item BIOS/UEFI firmware
  \item router firmware
  \item embedded-controller firmware
\end{itemize}

\medskip

\subsection{What is an operating system?}

An operating system is system software that manages computer hardware and provides services and abstractions to applications.

Examples:

\begin{itemize}
  \item Linux
  \item Windows
  \item macOS
  \item Android
\end{itemize}

\medskip

\subsection{What is a kernel?}

The kernel is the core component of an operating system responsible for managing resources such as:

\begin{itemize}
  \item CPU
  \item memory
  \item processes
  \item devices
  \item files
\end{itemize}

\medskip

\subsection{What is a process?}

A process is a \textbf{program in execution}.

A program is passive code; a process is the active execution instance of that code.

\medskip

\subsection{What is a thread?}

A thread is the smallest commonly scheduled unit of execution within a process.

Multiple threads of the same process typically share:

\begin{itemize}
  \item code
  \item heap
  \item global data
\end{itemize}

but have their own:

\begin{itemize}
  \item stack
  \item program counter
  \item registers/state
\end{itemize}

\medskip

\subsection{Process vs program?}

\begin{center}
\setlength{\tabcolsep}{6pt}
\renewcommand{\arraystretch}{1.2}
\begin{tabular}{|l|l|}
\hline
\textbf{Program} & \textbf{Process} \\
\hline
Passive & Active \\
Stored instructions & Executing instructions \\
No execution state & Has execution state \\
Can exist on disk & Exists while executing \\
\hline
\end{tabular}
\end{center}

\medskip

\subsection{What is a CPU?}

CPU stands for \textbf{Central Processing Unit}.

It executes instructions and performs arithmetic, logical and control operations.

Major conceptual components:

\begin{itemize}
  \item ALU
  \item Control Unit
  \item Registers
\end{itemize}

\medskip

\subsection{What is ALU?}

\textbf{Arithmetic Logic Unit} performs:

\begin{itemize}
  \item arithmetic operations
  \item logical operations
  \item comparisons
  \item bitwise operations
\end{itemize}

\medskip

\subsection{What is the Control Unit?}

The Control Unit coordinates CPU operations by fetching, decoding and controlling execution of instructions.

\medskip

\subsection{What is a register?}

A register is a very small, extremely fast storage location inside the CPU used to hold data, addresses or intermediate results.

\medskip

\subsection{Why is RAM faster than SSD?}

RAM is directly accessible by the memory subsystem and is designed for very low-latency random access.

SSD uses persistent flash storage and involves additional controller and interface operations.

\medskip

\subsection{Why is cache faster than RAM?}

Cache is implemented using faster memory technology, typically SRAM, and is physically close to or integrated with the CPU.

The hierarchy is approximately:

\begin{lstlisting}
Registers
   $\downarrow$
L1 Cache
   $\downarrow$
L2 Cache
   $\downarrow$
L3 Cache
   $\downarrow$
RAM
   $\downarrow$
SSD/HDD
\end{lstlisting}

As we go downward:

\begin{itemize}
  \item capacity generally increases
  \item cost per bit decreases
  \item latency increases
\end{itemize}

\medskip

\subsection{What is cache memory?}

Cache is a small, high-speed memory that stores frequently or recently used data/instructions to reduce access time.

\medskip

\subsection{What is locality of reference?}

Programs tend to access the same or nearby memory locations repeatedly.

Two important types:

\subsubsection{Temporal locality}

Recently accessed data is likely to be accessed again.

\subsubsection{Spatial locality}

Nearby memory locations are likely to be accessed soon.

Cache exploits both.

\medskip

\subsection{What is ROM?}

ROM means \textbf{Read Only Memory}.

It is non-volatile memory traditionally used to store firmware.

Modern ROM-like technologies include:

\begin{itemize}
  \item PROM
  \item EPROM
  \item EEPROM
  \item Flash
\end{itemize}

\medskip

\subsection{RAM vs ROM?}

\begin{center}
\setlength{\tabcolsep}{6pt}
\renewcommand{\arraystretch}{1.2}
\begin{tabular}{|l|l|}
\hline
\textbf{RAM} & \textbf{ROM} \\
\hline
Volatile & Non-volatile \\
Read/write & Primarily firmware/read-oriented \\
Working memory & Firmware/storage \\
Loses contents without power & Retains contents \\
\hline
\end{tabular}
\end{center}

\medskip

\subsection{What is volatile memory?}

Memory whose contents are lost when power is removed.

Example:

\begin{quote}
RAM\\
\end{quote}

\medskip

\subsection{What is non-volatile memory?}

Memory that retains information without continuous power.

Examples:

\begin{itemize}
  \item SSD
  \item HDD
  \item ROM
  \item flash memory
\end{itemize}

\medskip

\subsection{What is a bit?}

A bit is the smallest unit of digital information.

It can be:

\begin{lstlisting}
0 or 1
\end{lstlisting}

\medskip

\subsection{What is a byte?}

A byte consists of \textbf{8 bits}.

Therefore:

\begin{lstlisting}
1 byte = 8 bits
\end{lstlisting}

\medskip

\subsection{What is ASCII?}

ASCII is a character encoding standard originally using 7 bits to represent 128 characters.

For example:

\begin{lstlisting}
'A' = 65
'a' = 97
'0' = 48
\end{lstlisting}

\medskip

\subsection{What is Unicode?}

Unicode is a universal character-set standard designed to represent characters from languages and writing systems around the world.

UTF-8 is a popular Unicode encoding.

\medskip

\subsection{ASCII vs Unicode?}

ASCII represents a relatively small character set.

Unicode supports a vastly larger range of characters.

UTF-8 is backward-compatible with ASCII for ASCII characters.

\medskip

\subsection{What is binary?}

Binary is a base-2 number system using:

\begin{lstlisting}
0 and 1
\end{lstlisting}

Computers use binary because digital electronic circuits naturally represent two stable states.

\medskip

\subsection{Why do computers use binary?}

Digital hardware can reliably distinguish two states, such as:

\begin{lstlisting}
High voltage $\rightarrow$ 1
Low voltage  $\rightarrow$ 0
\end{lstlisting}

Binary therefore provides a robust representation for digital logic.

\medskip

\section{Number Systems \& Digital Logic}

\subsection{What number systems are commonly used in computing?}

\begin{itemize}
  \item Binary --- base 2
  \item Octal --- base 8
  \item Decimal --- base 10
  \item Hexadecimal --- base 16
\end{itemize}

\medskip

\subsection{Why is hexadecimal useful?}

Binary numbers become very long.

One hexadecimal digit represents exactly four binary bits.

Example:

\begin{lstlisting}
1111 = F
1010 = A
\end{lstlisting}

\medskip

\subsection{What is a logic gate?}

A logic gate is a digital circuit that performs a Boolean operation.

Common gates:

\begin{itemize}
  \item AND
  \item OR
  \item NOT
  \item NAND
  \item NOR
  \item XOR
  \item XNOR
\end{itemize}

\medskip

\subsection{Why are NAND and NOR called universal gates?}

Any Boolean function can be constructed using only NAND gates or only NOR gates.

\medskip

\subsection{AND vs OR?}

AND outputs 1 only when \textbf{all inputs are 1}.

OR outputs 1 when \textbf{at least one input is 1}.

\medskip

\subsection{XOR?}

XOR outputs 1 when the inputs are different.

For two inputs:

\begin{lstlisting}
0 XOR 0 = 0
0 XOR 1 = 1
1 XOR 0 = 1
1 XOR 1 = 0
\end{lstlisting}

\medskip

\subsection{Where is XOR useful?}

Common applications include:

\begin{itemize}
  \item parity checking
  \item bit toggling
  \item detecting differences
  \item cryptographic constructions
  \item finding a unique element using XOR properties
\end{itemize}

\medskip

\subsection{What is a flip-flop?}

A flip-flop is a sequential digital circuit capable of storing one bit of state.

\medskip

\subsection{Combinational vs sequential circuit?}

\textbf{Combinational:} output depends only on current input.

Examples:

\begin{itemize}
  \item adder
  \item multiplexer
  \item decoder
\end{itemize}

\textbf{Sequential:} output depends on current input and previous state.

Examples:

\begin{itemize}
  \item flip-flop
  \item register
  \item counter
\end{itemize}

\medskip

\section{Data Structures}

\subsection{What is a data structure?}

A data structure is a way of organizing and storing data so that operations can be performed efficiently.

\medskip

\subsection{Why do we need data structures?}

Because different operations require different performance characteristics.

For example:

\begin{itemize}
  \item Array $\rightarrow$ fast indexed access
  \item Hash table $\rightarrow$ fast average lookup
  \item Linked list $\rightarrow$ efficient insertion/deletion with a known node
  \item Heap $\rightarrow$ efficient priority operations
\end{itemize}

\medskip

\subsection{What are linear data structures?}

Elements are conceptually arranged sequentially.

Examples:

\begin{itemize}
  \item array
  \item linked list
  \item stack
  \item queue
\end{itemize}

\medskip

\subsection{What are non-linear data structures?}

Elements form hierarchical or graph-like relationships.

Examples:

\begin{itemize}
  \item tree
  \item graph
  \item heap
\end{itemize}

\medskip

\section{Arrays}

\subsection{What is an array?}

An array is a collection of elements stored in contiguous memory locations, typically with the same data type.

\medskip

\subsection{Why is array access O(1)?}

Because the address of an element can be calculated directly:

\begin{lstlisting}
address = base + index $\times$ element_size
\end{lstlisting}

Therefore no traversal is necessary.

\medskip

\subsection{Array advantages?}

\begin{itemize}
  \item O(1) indexed access
  \item cache-friendly
  \item simple structure
  \item low per-element overhead
\end{itemize}

\medskip

\subsection{Array disadvantages?}

\begin{itemize}
  \item fixed size in many traditional implementations
  \item insertion/deletion can require shifting
  \item contiguous memory requirement
\end{itemize}

\medskip

\subsection{Array insertion complexity?}

At the beginning:

\begin{lstlisting}
O(n)
\end{lstlisting}

At the end:

\begin{lstlisting}
O(1)
\end{lstlisting}

amortized for a dynamic array.

At an arbitrary position:

\begin{lstlisting}
O(n)
\end{lstlisting}

because elements may need shifting.

\medskip

\section{Linked Lists}

\subsection{What is a linked list?}

A linked list is a sequence of nodes where each node contains data and one or more links to other nodes.

\medskip

\subsection{Array vs linked list?}

\begin{center}
\setlength{\tabcolsep}{6pt}
\renewcommand{\arraystretch}{1.2}
\begin{tabular}{|l|l|}
\hline
\textbf{Array} & \textbf{Linked List} \\
\hline
Contiguous memory & Non-contiguous \\
O(1) indexing & O(n) indexing \\
Better cache locality & Poorer cache locality \\
Insertion can require shifting & Easy insertion with known node \\
Less pointer overhead & Pointer overhead \\
\hline
\end{tabular}
\end{center}

\medskip

\subsection{Types of linked lists?}

\begin{itemize}
  \item Singly linked list
  \item Doubly linked list
  \item Circular singly linked list
  \item Circular doubly linked list
\end{itemize}

\medskip

\subsection{What is a doubly linked list?}

Each node has links to both:

\begin{lstlisting}
previous node
next node
\end{lstlisting}

This allows traversal in both directions.

\medskip

\subsection{What is a circular linked list?}

The final node points back to the first node.

\medskip

\section{Stack}

\subsection{What is a stack?}

A stack is a linear data structure following:

\begin{quote}
\textbf{LIFO --- Last In, First Out}\\
\end{quote}

\medskip

\subsection{Stack applications?}

\begin{itemize}
  \item function call stack
  \item recursion
  \item undo operations
  \item expression evaluation
  \item parentheses matching
  \item DFS
\end{itemize}

\medskip

\subsection{What are stack operations?}

\begin{itemize}
  \item push
  \item pop
  \item peek/top
\end{itemize}

Usually O(1).

\medskip

\section{Queue}

\subsection{What is a queue?}

A queue follows:

\begin{quote}
\textbf{FIFO --- First In, First Out}\\
\end{quote}

\medskip

\subsection{Queue applications?}

\begin{itemize}
  \item CPU scheduling
  \item printer queues
  \item network buffering
  \item BFS
  \item request processing
\end{itemize}

\medskip

\subsection{What is a circular queue?}

A queue where the end wraps around to the beginning, allowing efficient reuse of freed positions.

\medskip

\subsection{Stack vs queue?}

\begin{lstlisting}
Stack:
Last in $\rightarrow$ First out

Queue:
First in $\rightarrow$ First out
\end{lstlisting}

\medskip

\section{Hash Table}

\subsection{What is a hash table?}

A hash table stores key-value pairs using a hash function to map keys to storage positions.

\medskip

\subsection{Average lookup complexity?}

Typically:

\begin{lstlisting}
O(1)
\end{lstlisting}

Average case.

Worst case can be:

\begin{lstlisting}
O(n)
\end{lstlisting}

depending on implementation and collisions.

\medskip

\subsection{What is a hash collision?}

A collision occurs when two different keys map to the same hash-table position.

\medskip

\subsection{How can collisions be handled?}

Common methods:

\begin{enumerate}
  \item Chaining
  \item Open addressing
\end{enumerate}

Open addressing includes:

\begin{itemize}
  \item linear probing
  \item quadratic probing
  \item double hashing
\end{itemize}

\medskip

\subsection{What is a hash function?}

A function that maps a key to a hash value/index.

A good hash function should distribute keys approximately uniformly.

\medskip

\section{Algorithms \& Complexity}

\subsection{What is an algorithm?}

An algorithm is a finite, well-defined sequence of steps for solving a problem.

\medskip

\subsection{What makes a good algorithm?}

It should ideally have:

\begin{itemize}
  \item correctness
  \item termination
  \item efficiency
  \item clarity
  \item appropriate resource usage
\end{itemize}

\medskip

\subsection{What is time complexity?}

Time complexity describes how an algorithm's computational work grows as input size increases.

\medskip

\subsection{What is space complexity?}

Space complexity describes how additional memory requirements grow with input size.

\medskip

\subsection{What is Big-O?}

Big-O describes an asymptotic upper bound on growth rate.

Common complexities:

\begin{lstlisting}
O(1)
O(log n)
O(n)
O(n log n)
O(n\textsuperscript{2})
O(2)
O(n!)
\end{lstlisting}

\medskip

\subsection{Which is better: O(n) or O(n\textsuperscript{2})?}

Generally O(n), especially as n becomes large.

\medskip

\subsection{What is O(1)?}

Constant time.

The operation does not grow with input size.

Example:

\begin{lstlisting}
array[index]
\end{lstlisting}

\medskip

\subsection{What is O(log n)?}

The problem size is reduced by a constant factor at each step.

Classic example:

\begin{quote}
Binary search.\\
\end{quote}

\medskip

\subsection{Why is binary search O(log n)?}

Because each comparison approximately halves the search space.

For n elements:

\begin{lstlisting}
n
n/2
n/4
n/8
...
1
\end{lstlisting}

The number of steps is approximately:

\begin{lstlisting}
log n
\end{lstlisting}

\medskip

\subsection{What is O(n log n)?}

A common complexity for efficient comparison sorting algorithms.

Examples:

\begin{itemize}
  \item Merge sort
  \item Heap sort
  \item average-case quicksort
\end{itemize}

\medskip

\section{Searching}

\subsection{Linear search?}

Checks elements sequentially.

Complexity:

\begin{lstlisting}
Best:    O(1)
Average: O(n)
Worst:   O(n)
\end{lstlisting}

\medskip

\subsection{Binary search?}

Binary search repeatedly divides a \textbf{sorted} search space in half.

Complexity:

\begin{lstlisting}
O(log n)
\end{lstlisting}

\medskip

\subsection{Why must binary search normally use sorted data?}

Because it decides which half to discard based on ordering.

Without ordering, discarding half would not be justified.

\medskip

\section{Sorting}

\subsection{What is sorting?}

Arranging data according to a specified order.

Example:

\begin{lstlisting}
5 2 9 1
$\downarrow$
1 2 5 9
\end{lstlisting}

\medskip

\subsection{Bubble sort complexity?}

Worst case:

\begin{lstlisting}
O(n\textsuperscript{2})
\end{lstlisting}

Best case with optimization:

\begin{lstlisting}
O(n)
\end{lstlisting}

\medskip

\subsection{Selection sort complexity?}

Generally:

\begin{lstlisting}
O(n\textsuperscript{2})
\end{lstlisting}

even in the best case.

\medskip

\subsection{Insertion sort complexity?}

\begin{lstlisting}
Best:  O(n)
Worst: O(n\textsuperscript{2})
\end{lstlisting}

It works particularly well for small or nearly sorted datasets.

\medskip

\subsection{Merge sort complexity?}

\begin{lstlisting}
Time: O(n log n)
Space: O(n)
\end{lstlisting}

It uses divide and conquer.

\medskip

\subsection{Quick sort complexity?}

Average:

\begin{lstlisting}
O(n log n)
\end{lstlisting}

Worst:

\begin{lstlisting}
O(n\textsuperscript{2})
\end{lstlisting}

depending on pivot choices.

\medskip

\subsection{Heap sort complexity?}

\begin{lstlisting}
O(n log n)
\end{lstlisting}

with O(1) auxiliary array space for a typical in-place implementation.

\medskip

\subsection{Which sorting algorithm is stable?}

Examples:

\begin{itemize}
  \item Merge sort
  \item Insertion sort
  \item Bubble sort
\end{itemize}

Quicksort and heapsort are not inherently stable in their common in-place forms.

\medskip

\subsection{What is a stable sort?}

If two elements have equal keys, their relative order remains unchanged.

\medskip

\section{Recursion}

\subsection{What is recursion?}

Recursion occurs when a function directly or indirectly calls itself.

A recursive solution normally needs:

\begin{enumerate}
  \item Base case
  \item Recursive case
\end{enumerate}

\medskip

\subsection{What happens if there is no base case?}

The recursion may continue indefinitely until resources are exhausted, typically causing stack overflow.

\medskip

\subsection{Recursion vs iteration?}

Recursion uses function calls and stack frames.

Iteration uses loops.

Recursion can make tree and divide-and-conquer algorithms easier to express, but may consume additional stack memory.

\medskip

\section{Object-Oriented Programming}

\subsection{What is OOP?}

Object-Oriented Programming organizes software around objects that combine state and behavior.

\medskip

\subsection{Four pillars of OOP?}

\begin{enumerate}
  \item Encapsulation
  \item Abstraction
  \item Inheritance
  \item Polymorphism
\end{enumerate}

\medskip

\subsection{What is a class?}

A class is a blueprint or user-defined type describing data and behavior.

\medskip

\subsection{What is an object?}

An object is an instance of a class.

\medskip

\subsection{Class vs object?}

\begin{lstlisting}
Class  = blueprint
Object = actual instance
\end{lstlisting}

Example:

\begin{lstlisting}
Class: Car
Objects: Toyota, BMW, Honda
\end{lstlisting}

\medskip

\subsection{What is encapsulation?}

Encapsulation combines data and related methods while controlling access to internal state.

For example:

\begin{lstlisting}
private double balance;

public double getBalance() {
    return balance;
}
\end{lstlisting}

\medskip

\subsection{Why use private variables?}

To prevent uncontrolled external modification and preserve invariants.

\medskip

\subsection{What is abstraction?}

Abstraction exposes essential behavior while hiding unnecessary implementation details.

Example:

\begin{lstlisting}
ATM user:
withdraw()
deposit()
\end{lstlisting}

The user does not need to know the internal banking implementation.

\medskip

\subsection{Encapsulation vs abstraction?}

\textbf{Encapsulation:} controls access and bundles state with behavior.

\textbf{Abstraction:} hides unnecessary implementation complexity and exposes essential concepts.

\medskip

\subsection{What is inheritance?}

Inheritance allows a class to derive properties and behavior from another class.

Example:

\begin{lstlisting}
Vehicle
   $\downarrow$
Car
\end{lstlisting}

\medskip

\subsection{Why use inheritance?}

For:

\begin{itemize}
  \item code reuse
  \item specialization
  \item hierarchical modeling
  \item polymorphism
\end{itemize}

However, inheritance should not be used merely for code reuse when composition is more appropriate.

\medskip

\subsection{What is polymorphism?}

Polymorphism means one interface can represent different underlying implementations.

Two common forms:

\begin{itemize}
  \item compile-time polymorphism
  \item runtime polymorphism
\end{itemize}

\medskip

\subsection{Method overloading?}

Same method name but different parameter lists.

This is compile-time polymorphism.

\medskip

\subsection{Method overriding?}

A subclass provides a new implementation of an inherited method.

This supports runtime polymorphism when dynamic dispatch is used.

\medskip

\subsection{Overloading vs overriding?}

\begin{center}
\setlength{\tabcolsep}{6pt}
\renewcommand{\arraystretch}{1.2}
\begin{tabular}{|l|l|}
\hline
\textbf{Overloading} & \textbf{Overriding} \\
\hline
Same class commonly & Parent-child relationship \\
Different parameters & Same method signature \\
Compile-time selection & Runtime dispatch \\
Different behavior signatures & Replaces inherited implementation \\
\hline
\end{tabular}
\end{center}

\medskip

\subsection{What is a constructor?}

A constructor initializes an object when it is created.

\medskip

\subsection{Does a constructor have a return type?}

No.

Not even \texttt{void}.

\medskip

\subsection{What is a destructor?}

A destructor is a special mechanism used in languages such as C++ to perform cleanup when an object's lifetime ends.

Java does not use C++-style destructors.

\medskip

\subsection{What is an interface?}

An interface defines a contract that implementing classes must satisfy.

In Java, interfaces can also contain default and static methods.

\medskip

\subsection{Abstract class vs interface?}

An abstract class can contain:

\begin{itemize}
  \item state
  \item constructors
  \item concrete methods
  \item abstract methods
\end{itemize}

An interface primarily defines a contract and supports multiple interface inheritance in Java.

\medskip

\subsection{What is composition?}

Composition means building a class using objects of other classes.

Example:

\begin{lstlisting}
Car has an Engine
\end{lstlisting}

This is a \textbf{has-a} relationship.

\medskip

\subsection{Inheritance vs composition?}

Inheritance:

\begin{quote}
\textbf{is-a}\\
\end{quote}

Composition:

\begin{quote}
\textbf{has-a}\\
\end{quote}

In many designs, composition provides looser coupling and greater flexibility.

\medskip

\section{C / C++}

\subsection{What is C?}

C is a procedural, compiled, general-purpose programming language known for low-level memory access and efficiency.

\medskip

\subsection{What is C++?}

C++ is a multi-paradigm language supporting procedural, object-oriented, generic and other programming styles.

\medskip

\subsection{C vs C++?}

C primarily emphasizes procedural programming.

C++ adds features such as:

\begin{itemize}
  \item classes
  \item inheritance
  \item polymorphism
  \item templates
  \item exceptions
  \item STL
\end{itemize}

\medskip

\subsection{What is a pointer?}

A pointer is a variable that stores the memory address of another object or function, depending on the pointer type.

\medskip

\subsection{What is a reference in C++?}

A reference is an alias for an existing object.

Example:

\begin{lstlisting}
int x = 10;
int &r = x;
\end{lstlisting}

\medskip

\subsection{Pointer vs reference?}

A pointer:

\begin{itemize}
  \item stores an address
  \item can be null
  \item can be reassigned
\end{itemize}

A reference:

\begin{itemize}
  \item aliases an existing object
  \item normally must be initialized
  \item cannot be reseated to refer to another object
\end{itemize}

\medskip

\subsection{What is a null pointer?}

A pointer that does not point to a valid object.

In modern C++:

\begin{lstlisting}
nullptr
\end{lstlisting}

is preferred.

\medskip

\subsection{What is a dangling pointer?}

A pointer that refers to memory whose lifetime has ended.

\medskip

\subsection{What is a memory leak?}

Allocated memory that is no longer reachable and therefore cannot be properly released.

\medskip

\subsection{Stack vs heap memory?}

\subsubsection{Stack}

Usually stores:

\begin{itemize}
  \item local variables
  \item function call frames
\end{itemize}

Fast and automatically managed.

\subsubsection{Heap}

Used for dynamically allocated objects.

Lifetime can extend beyond the current function, depending on allocation and ownership.

\medskip

\subsection{What is dynamic memory allocation?}

Allocating memory during program execution.

C:

\begin{lstlisting}
malloc()
calloc()
realloc()
free()
\end{lstlisting}

C++:

\begin{lstlisting}
new
delete
\end{lstlisting}

Modern C++ generally prefers RAII and smart pointers over raw \texttt{new/delete}.

\medskip

\subsection{malloc vs calloc?}

\texttt{malloc} allocates uninitialized memory.

\texttt{calloc} allocates memory and initializes the allocated bytes to zero.

\medskip

\subsection{malloc vs new?}

\texttt{malloc} is a C library allocation function and returns raw storage.

\texttt{new} is a C++ operator that allocates storage and constructs an object.

\medskip

\subsection{What is free()?}

It releases memory allocated using C allocation functions such as \texttt{malloc}.

\medskip

\subsection{What is delete?}

It destroys an object allocated with \texttt{new}.

\medskip

\subsection{Why should we not mix malloc/free with new/delete?}

Because they use different allocation/deallocation mechanisms.

Correct pairs:

\begin{lstlisting}
malloc  $\rightarrow$ free
new     $\rightarrow$ delete
new[]   $\rightarrow$ delete[]
\end{lstlisting}

\medskip

\subsection{What is a segmentation fault?}

An invalid memory access that causes the operating system to terminate the process.

Common causes:

\begin{itemize}
  \item dereferencing invalid pointers
  \item accessing outside an array
  \item using freed memory
\end{itemize}

\medskip

\subsection{What is a structure?}

A structure groups related variables, potentially of different types, under one user-defined type.

\medskip

\subsection{Structure vs class in C++?}

The main default access difference:

\begin{lstlisting}
struct $\rightarrow$ public
class  $\rightarrow$ private
\end{lstlisting}

Otherwise, C++ supports many of the same capabilities in both.

\medskip

\subsection{What is STL?}

STL means \textbf{Standard Template Library}.

It provides reusable:

\begin{itemize}
  \item containers
  \item algorithms
  \item iterators
  \item function objects/utilities
\end{itemize}

Examples:

\begin{lstlisting}
vector
map
set
queue
stack
sort()
\end{lstlisting}

\medskip

\subsection{vector vs array?}

\texttt{std::array} has fixed size.

\texttt{std::vector} is dynamically resizable.

Both provide O(1) random access.

\medskip

\subsection{vector vs linked list?}

\texttt{vector} usually has much better cache locality and O(1) random access.

\texttt{list} provides constant-time insertion/deletion at a known iterator position but O(n) traversal and significant pointer overhead.

\medskip

\subsection{map vs unordered\_map?}

\texttt{map} is typically implemented as a balanced tree:

\begin{lstlisting}
lookup: O(log n)
ordered
\end{lstlisting}

\texttt{unordered\_map} is hash-based:

\begin{lstlisting}
average lookup: O(1)
unordered
\end{lstlisting}

\medskip

\subsection{set vs unordered\_set?}

\texttt{set} maintains ordered unique elements.

\texttt{unordered\_set} stores unique elements using hashing without ordering.

\medskip

\section{Java}

\subsection{Why is Java platform independent?}

Java source is compiled into bytecode, which runs on a Java Virtual Machine.

\begin{lstlisting}
Java source
    $\downarrow$
Bytecode
    $\downarrow$
JVM
    $\downarrow$
Operating system
\end{lstlisting}

\medskip

\subsection{What is JVM?}

\textbf{Java Virtual Machine} executes Java bytecode.

\medskip

\subsection{JDK vs JRE vs JVM?}

\subsubsection{JVM}

Executes bytecode.

\subsubsection{JRE}

JVM + libraries/runtime components.

\subsubsection{JDK}

JRE/runtime + development tools such as the compiler.

\medskip

\subsection{Why is Java not purely object-oriented?}

Because it has primitive types such as:

\begin{lstlisting}
int
char
boolean
double
\end{lstlisting}

which are not objects.

\medskip

\subsection{What is garbage collection?}

Automatic identification and reclamation of heap objects that are no longer reachable.

\medskip

\subsection{Can Java have memory leaks?}

Yes.

Garbage collection prevents many manual-memory errors, but objects can remain reachable unnecessarily.

Example:

\begin{lstlisting}
Static collection
     $\downarrow$
Object retained forever
\end{lstlisting}

\medskip

\subsection{What is final?}

Depending on context:

\begin{itemize}
  \item final variable $\rightarrow$ cannot be reassigned
  \item final method $\rightarrow$ cannot be overridden
  \item final class $\rightarrow$ cannot be extended
\end{itemize}

\medskip

\subsection{What is static?}

A static member belongs to the class rather than to individual instances.

\medskip

\subsection{What is an exception?}

An exception is an event representing an abnormal condition during program execution.

\medskip

\subsection{Checked vs unchecked exception?}

Checked exceptions are checked by the compiler.

Unchecked exceptions generally derive from \texttt{RuntimeException}.

\medskip

\section{Operating Systems}

\subsection{Main functions of an OS?}

\begin{itemize}
  \item process management
  \item memory management
  \item file management
  \item device management
  \item security/protection
  \item networking
  \item resource allocation
\end{itemize}

\medskip

\subsection{What is process scheduling?}

Selecting which ready process/thread should receive CPU time.

\medskip

\subsection{Common CPU scheduling algorithms?}

\begin{itemize}
  \item FCFS
  \item SJF
  \item SRTF
  \item Round Robin
  \item Priority Scheduling
  \item Multilevel Queue
  \item Multilevel Feedback Queue
\end{itemize}

\medskip

\subsection{What is FCFS?}

\textbf{First Come First Served.}

Processes execute in arrival order.

Simple, but can suffer from the \textbf{convoy effect}.

\medskip

\subsection{What is SJF?}

\textbf{Shortest Job First.}

The process with the smallest estimated CPU burst is selected.

It minimizes average waiting time under ideal assumptions, but requires burst-time estimation and can starve long jobs.

\medskip

\subsection{What is Round Robin?}

Each process receives a fixed time slice called a \textbf{quantum}.

It is widely suited to interactive/time-sharing systems.

\medskip

\subsection{What happens if the time quantum is too small?}

There are too many context switches, increasing overhead.

\medskip

\subsection{What if quantum is too large?}

Round Robin approaches FCFS behavior.

\medskip

\subsection{What is context switching?}

Saving the state of one running task and restoring another task's state so the CPU can switch execution.

\medskip

\subsection{What is a PCB?}

\textbf{Process Control Block.}

It stores process information such as:

\begin{itemize}
  \item process ID
  \item state
  \item program counter
  \item CPU registers
  \item scheduling information
  \item memory-management information
\end{itemize}

\medskip

\section{Threads}

\subsection{Process vs thread?}

A process has its own address space.

Threads within a process generally share the process's address space and resources.

Threads are therefore cheaper to create/switch than separate processes in many systems.

\medskip

\subsection{What is multithreading?}

Executing multiple threads concurrently within a process.

\medskip

\subsection{Advantage of multithreading?}

\begin{itemize}
  \item responsiveness
  \item resource sharing
  \item concurrency
  \item parallelism on multicore systems
\end{itemize}

\medskip

\subsection{Concurrency vs parallelism?}

\textbf{Concurrency:} multiple tasks make progress during overlapping periods.

\textbf{Parallelism:} multiple tasks literally execute simultaneously on multiple processing resources.

\medskip

\section{Synchronization}

\subsection{What is a race condition?}

A race condition occurs when program behavior depends on the timing/interleaving of concurrent operations.

\medskip

\subsection{What is a critical section?}

A section of code that accesses shared resources and therefore must be protected from unsafe concurrent access.

\medskip

\subsection{What is a mutex?}

A mutex provides mutual exclusion so that only one thread can hold the lock at a time.

\medskip

\subsection{What is a semaphore?}

A semaphore is a synchronization primitive containing a counter used to control access to resources.

Types include:

\begin{itemize}
  \item binary semaphore
  \item counting semaphore
\end{itemize}

\medskip

\subsection{Mutex vs semaphore?}

A mutex is primarily an ownership-based mutual-exclusion mechanism.

A semaphore is a signaling/resource-counting mechanism and does not necessarily have ownership semantics.

\medskip

\section{Deadlock}

\subsection{What is deadlock?}

A situation where processes/threads wait indefinitely for resources held by one another.

\medskip

\subsection{Four necessary conditions for deadlock?}

\textbf{Coffman conditions:}

\begin{enumerate}
  \item Mutual exclusion
  \item Hold and wait
  \item No preemption
  \item Circular wait
\end{enumerate}

\medskip

\subsection{How can deadlock be prevented?}

Break at least one necessary condition.

For example, impose a strict resource acquisition ordering to prevent circular wait.

\medskip

\subsection{Deadlock vs starvation?}

\textbf{Deadlock:} involved processes wait indefinitely for one another.

\textbf{Starvation:} a process waits indefinitely because resources/CPU repeatedly go to others.

\medskip

\subsection{What is livelock?}

Processes remain active and repeatedly change state in response to one another but make no useful progress.

\medskip

\section{Memory Management}

\subsection{What is virtual memory?}

Virtual memory allows processes to use an address space larger than physical RAM by using mechanisms such as paging and secondary storage.

\medskip

\subsection{What is paging?}

Paging divides:

\begin{itemize}
  \item virtual memory into pages
  \item physical memory into frames
\end{itemize}

Pages can be mapped to frames.

\medskip

\subsection{What is a page fault?}

A page fault occurs when a process accesses a virtual page that is not currently available in physical memory.

The OS then handles the fault, potentially loading the page from storage.

\medskip

\subsection{What is segmentation?}

Memory is divided into logical variable-sized segments such as:

\begin{itemize}
  \item code
  \item data
  \item stack
\end{itemize}

Modern systems commonly rely heavily on paging, though segmentation concepts remain important historically and architecturally.

\medskip

\subsection{What is fragmentation?}

Two types:

\subsubsection{Internal fragmentation}

Allocated block contains unused space.

\subsubsection{External fragmentation}

Free memory exists but is split into small non-contiguous pieces.

\medskip

\section{DBMS}

\subsection{What is a DBMS?}

A Database Management System is software used to store, organize, retrieve and manage data.

Examples:

\begin{itemize}
  \item PostgreSQL
  \item MySQL
  \item Oracle Database
  \item SQL Server
\end{itemize}

\medskip

\subsection{What is a database?}

An organized collection of data that can be efficiently stored and accessed.

\medskip

\subsection{DBMS vs RDBMS?}

An RDBMS organizes data primarily into related tables and generally follows the relational model.

Examples:

\begin{itemize}
  \item MySQL
  \item PostgreSQL
  \item Oracle
\end{itemize}

\medskip

\subsection{What is a table?}

A relational table represents data using rows and columns.

\medskip

\subsection{What is a row?}

A row represents one record/tuple.

\medskip

\subsection{What is a column?}

A column represents an attribute/field of the records.

\medskip

\section{Keys}

\subsection{What is a primary key?}

A primary key uniquely identifies each row in a table.

Properties:

\begin{itemize}
  \item unique
  \item not null
\end{itemize}

\medskip

\subsection{What is a foreign key?}

A foreign key references a key in another table and establishes a relationship between tables.

\medskip

\subsection{Primary key vs unique key?}

A primary key uniquely identifies rows and cannot be NULL.

A table can have multiple UNIQUE constraints, and NULL behavior depends on the DBMS.

\medskip

\subsection{What is a candidate key?}

A minimal set of attributes capable of uniquely identifying a tuple.

\medskip

\subsection{What is a super key?}

Any set of attributes that uniquely identifies a tuple.

A candidate key is a \textbf{minimal super key}.

\medskip

\section{Normalization}

\subsection{What is normalization?}

Normalization is the process of organizing relational data to reduce redundancy and undesirable dependencies.

\medskip

\subsection{What is 1NF?}

A relation is in First Normal Form when attributes contain atomic values and there are no repeating groups in the relational sense.

\medskip

\subsection{What is 2NF?}

2NF requires:

\begin{itemize}
  \item 1NF
  \item no partial dependency of a non-key attribute on part of a composite candidate key
\end{itemize}

\medskip

\subsection{What is 3NF?}

3NF eliminates relevant transitive dependencies of non-key attributes on candidate keys.

A common simplified formulation:

\begin{quote}
Every non-key attribute depends on the key, the whole key, and nothing but the key.\\
\end{quote}

\medskip

\subsection{What is BCNF?}

A relation is in BCNF if every determinant is a candidate key.

BCNF is stricter than 3NF.

\medskip

\subsection{Why normalize databases?}

To reduce:

\begin{itemize}
  \item redundancy
  \item update anomalies
  \item insertion anomalies
  \item deletion anomalies
\end{itemize}

\medskip

\subsection{Can over-normalization be bad?}

Yes.

Highly normalized designs may require many joins and can hurt read performance.

Sometimes \textbf{denormalization} is intentionally used for performance.

\medskip

\section{SQL}

\subsection{What is SQL?}

SQL stands for \textbf{Structured Query Language} and is used to interact with relational databases.

\medskip

\subsection{SQL categories?}

Common classification:

\subsubsection{DDL}

\begin{lstlisting}
CREATE
ALTER
DROP
TRUNCATE
\end{lstlisting}

\subsubsection{DML}

\begin{lstlisting}
INSERT
UPDATE
DELETE
\end{lstlisting}

\subsubsection{DQL}

\begin{lstlisting}
SELECT
\end{lstlisting}

\subsubsection{DCL}

\begin{lstlisting}
GRANT
REVOKE
\end{lstlisting}

\subsubsection{TCL}

\begin{lstlisting}
COMMIT
ROLLBACK
SAVEPOINT
\end{lstlisting}

\medskip

\subsection{DELETE vs TRUNCATE vs DROP?}

\subsubsection{DELETE}

Removes rows and can use a WHERE condition.

\subsubsection{TRUNCATE}

Removes all rows efficiently, subject to DBMS-specific transactional behavior.

\subsubsection{DROP}

Removes the table/object itself.

\medskip

\subsection{What is a JOIN?}

A JOIN combines rows from multiple tables based on a related condition.

\medskip

\subsection{Types of joins?}

\begin{itemize}
  \item INNER JOIN
  \item LEFT JOIN
  \item RIGHT JOIN
  \item FULL OUTER JOIN
  \item CROSS JOIN
  \item SELF JOIN
\end{itemize}

\medskip

\subsection{INNER JOIN?}

Returns rows satisfying the join condition in both tables.

\medskip

\subsection{LEFT JOIN?}

Returns all rows from the left table and matching rows from the right table.

If no match exists, right-side columns become NULL.

\medskip

\subsection{What is a self join?}

Joining a table with itself.

Example:

\begin{lstlisting}
Employee $\rightarrow$ Manager
\end{lstlisting}

where both are stored in the Employee table.

\medskip

\subsection{WHERE vs HAVING?}

\texttt{WHERE} filters rows before grouping.

\texttt{HAVING} filters groups after \texttt{GROUP BY}.

\medskip

\subsection{GROUP BY?}

Groups rows according to specified columns, usually for aggregate calculations.

\medskip

\subsection{Common aggregate functions?}

\begin{lstlisting}
COUNT()
SUM()
AVG()
MIN()
MAX()
\end{lstlisting}

\medskip

\subsection{What is a subquery?}

A query nested inside another SQL statement.

\medskip

\subsection{JOIN vs subquery?}

Neither is universally better.

JOINs are often natural for combining related tables.

Subqueries can make certain existence/filtering logic clearer.

The optimizer may transform equivalent queries into similar execution plans.

\medskip

\subsection{What is a view?}

A view is a virtual table defined by a query.

It generally stores the query definition rather than a separate copy of the underlying data.

\medskip

\subsection{What is an index?}

An index is a data structure that speeds up data retrieval at the cost of additional storage and write overhead.

\medskip

\subsection{Why can too many indexes be harmful?}

Every insert/update/delete may need to maintain those indexes.

Therefore indexes improve reads but can increase write cost and storage usage.

\medskip

\section{Transactions}

\subsection{What is a transaction?}

A transaction is a logical unit of database work that should be completed consistently.

\medskip

\subsection{ACID properties?}

\subsubsection{Atomicity}

All or nothing.

\subsubsection{Consistency}

Preserves database rules/invariants.

\subsubsection{Isolation}

Concurrent transactions behave according to the isolation model.

\subsubsection{Durability}

Committed changes survive failures according to the DBMS's guarantees.

\medskip

\subsection{COMMIT?}

Makes transaction changes permanent.

\medskip

\subsection{ROLLBACK?}

Undoes uncommitted transaction changes.

\medskip

\subsection{What is a transaction isolation level?}

It determines how one transaction can observe changes made by concurrent transactions.

Common levels:

\begin{itemize}
  \item Read Uncommitted
  \item Read Committed
  \item Repeatable Read
  \item Serializable
\end{itemize}

Some DBMSs also implement snapshot/MVCC-based semantics.

\medskip

\section{Computer Networks}

\subsection{What is a computer network?}

A network is a collection of interconnected devices that communicate and share resources.

\medskip

\subsection{LAN?}

\textbf{Local Area Network.}

Covers a relatively small area such as:

\begin{itemize}
  \item office
  \item home
  \item campus
\end{itemize}

\medskip

\subsection{WAN?}

\textbf{Wide Area Network.}

Covers large geographic areas.

The Internet is the largest example of an interconnected collection of networks.

\medskip

\subsection{What is the Internet?}

The Internet is a global network of interconnected networks communicating primarily using the TCP/IP protocol suite.

\medskip

\section{OSI Model}

\subsection{Seven OSI layers?}

From bottom to top:

\begin{enumerate}
  \item Physical
  \item Data Link
  \item Network
  \item Transport
  \item Session
  \item Presentation
  \item Application
\end{enumerate}

Mnemonic:

\begin{quote}
\textbf{Please Do Not Throw Sausage Pizza Away}\\
\end{quote}

\medskip

\subsection{TCP/IP model?}

Common four-layer representation:

\begin{enumerate}
  \item Link
  \item Internet
  \item Transport
  \item Application
\end{enumerate}

Some textbooks use a five-layer model by separating physical and data-link layers.

\medskip

\subsection{What is the Physical layer?}

Deals with transmission of raw bits over physical media.

Examples:

\begin{itemize}
  \item cables
  \item radio
  \item connectors
  \item electrical/optical signals
\end{itemize}

\medskip

\subsection{Data Link layer?}

Provides node-to-node delivery over a local link.

Examples:

\begin{itemize}
  \item Ethernet
  \item Wi-Fi MAC
  \item switches operate primarily here
\end{itemize}

\medskip

\subsection{Network layer?}

Responsible for logical addressing and routing.

Main example:

\begin{quote}
IP\\
\end{quote}

\medskip

\subsection{Transport layer?}

Provides end-to-end transport services.

Major protocols:

\begin{itemize}
  \item TCP
  \item UDP
\end{itemize}

\medskip

\subsection{Application layer?}

Provides network services to applications.

Examples:

\begin{itemize}
  \item HTTP
  \item DNS
  \item SMTP
  \item FTP
  \item SSH
\end{itemize}

\medskip

\section{TCP / UDP}

\subsection{TCP?}

Transmission Control Protocol is:

\begin{itemize}
  \item connection-oriented
  \item reliable
  \item ordered
  \item flow-controlled
  \item congestion-controlled
\end{itemize}

\medskip

\subsection{UDP?}

User Datagram Protocol is connectionless and provides minimal transport functionality without TCP's reliability guarantees.

\medskip

\subsection{TCP vs UDP?}

\begin{center}
\setlength{\tabcolsep}{6pt}
\renewcommand{\arraystretch}{1.2}
\begin{tabular}{|l|l|}
\hline
\textbf{TCP} & \textbf{UDP} \\
\hline
Reliable & Best-effort \\
Ordered & No delivery ordering guarantee \\
Connection-oriented & Connectionless \\
More overhead & Lower overhead \\
Flow/congestion control & No TCP-style mechanisms \\
\hline
\end{tabular}
\end{center}

\medskip

\subsection{Where is UDP useful?}

\begin{itemize}
  \item DNS
  \item real-time media
  \item online gaming
  \item VoIP
  \item QUIC's underlying transport is UDP
\end{itemize}

Applications may implement their own reliability where needed.

\medskip

\section{IP}

\subsection{What is an IP address?}

A logical address used to identify an interface in an IP network.

\medskip

\subsection{IPv4?}

IPv4 uses \textbf{32-bit addresses}.

Example:

\begin{lstlisting}
192.168.1.10
\end{lstlisting}

\medskip

\subsection{IPv6?}

IPv6 uses \textbf{128-bit addresses}.

Example:

\begin{lstlisting}
2001:db8::1
\end{lstlisting}

\medskip

\subsection{Why was IPv6 introduced?}

Primarily because IPv4 has a limited address space.

IPv6 provides vastly more addresses and includes other architectural improvements.

\medskip

\subsection{Private IPv4 addresses?}

Common private ranges:

\begin{lstlisting}
10.0.0.0/8
172.16.0.0/12
192.168.0.0/16
\end{lstlisting}

\medskip

\subsection{What is NAT?}

Network Address Translation maps addresses between network domains, commonly translating private internal addresses to public addresses.

It also helps conserve IPv4 address space.

\medskip

\section{Subnetting}

\subsection{What is subnetting?}

Subnetting divides an IP network into smaller logical networks.

Benefits:

\begin{itemize}
  \item efficient address allocation
  \item isolation
  \item routing control
  \item broadcast-domain reduction
\end{itemize}

\medskip

\subsection{What does /24 mean?}

It means the first 24 bits are network-prefix bits.

IPv4 \texttt{/24} corresponds to:

\begin{lstlisting}
255.255.255.0
\end{lstlisting}

\medskip

\subsection{/27?}

\begin{lstlisting}
255.255.255.224
\end{lstlisting}

It provides:

\begin{lstlisting}
32 addresses
\end{lstlisting}

per subnet.

Traditionally:

\begin{lstlisting}
30 usable host addresses
\end{lstlisting}

for an ordinary IPv4 subnet excluding network and broadcast addresses.

\medskip

\subsection{What is a default gateway?}

The router/interface to which a host sends traffic destined for networks outside its local subnet.

\medskip

\section{Network Devices}

\subsection{Hub vs switch?}

A hub broadcasts incoming signals to all ports.

A switch forwards Ethernet frames based on MAC addresses.

\medskip

\subsection{Switch vs router?}

A switch primarily connects devices within a LAN.

A router forwards packets between different IP networks.

\medskip

\subsection{What is a MAC address?}

A MAC address is a link-layer identifier associated with a network interface.

Ethernet MAC addresses are typically 48 bits.

\medskip

\subsection{MAC vs IP?}

MAC:

\begin{quote}
Link-layer identification.\\
\end{quote}

IP:

\begin{quote}
Network-layer logical addressing/routing.\\
\end{quote}

\medskip

\section{DNS}

\subsection{What is DNS?}

Domain Name System translates domain names into information such as IP addresses.

Example:

\begin{lstlisting}
google.com
    $\downarrow$
IP address
\end{lstlisting}

\medskip

\subsection{Why is DNS called hierarchical?}

Because DNS is organized into levels:

\begin{lstlisting}
Root
 $\downarrow$
TLD
 $\downarrow$
Authoritative domain
 $\downarrow$
Subdomain/host
\end{lstlisting}

\medskip

\subsection{What is DHCP?}

Dynamic Host Configuration Protocol automatically provides network configuration such as:

\begin{itemize}
  \item IP address
  \item subnet mask/prefix
  \item default gateway
  \item DNS server
\end{itemize}

\medskip

\section{HTTP}

\subsection{What is HTTP?}

HyperText Transfer Protocol is an application-layer protocol used for communication between clients and web servers.

\medskip

\subsection{HTTP vs HTTPS?}

HTTPS is HTTP carried over a secure TLS connection.

It provides:

\begin{itemize}
  \item confidentiality
  \item integrity
  \item server authentication
\end{itemize}

when TLS is correctly configured.

\medskip

\subsection{Common HTTP methods?}

\begin{itemize}
  \item GET
  \item POST
  \item PUT
  \item PATCH
  \item DELETE
  \item HEAD
  \item OPTIONS
\end{itemize}

\medskip

\subsection{GET vs POST?}

GET generally retrieves a resource and is intended to be safe/idempotent.

POST generally submits data for processing or creates subordinate resources and is not inherently idempotent.

\medskip

\subsection{What is REST?}

REST is an architectural style for distributed systems.

RESTful APIs commonly use:

\begin{itemize}
  \item resources
  \item HTTP methods
  \item stateless requests
  \item representations such as JSON
\end{itemize}

\medskip

\subsection{What is an API?}

Application Programming Interface.

It defines how software components communicate.

\medskip

\section{Software Engineering}

\subsection{What is software engineering?}

The systematic application of engineering principles to the development, operation and maintenance of software.

\medskip

\subsection{SDLC?}

Software Development Life Cycle.

Common stages:

\begin{lstlisting}
Requirements
    $\downarrow$
Design
    $\downarrow$
Implementation
    $\downarrow$
Testing
    $\downarrow$
Deployment
    $\downarrow$
Maintenance
\end{lstlisting}

\medskip

\subsection{Waterfall model?}

A largely sequential development model where phases are performed in planned stages.

Useful when requirements are stable, but less flexible when requirements change frequently.

\medskip

\subsection{Agile?}

An iterative development approach emphasizing:

\begin{itemize}
  \item incremental delivery
  \item customer feedback
  \item adaptability
  \item collaboration
\end{itemize}

\medskip

\subsection{Scrum?}

An Agile framework using:

\begin{itemize}
  \item sprints
  \item product backlog
  \item sprint backlog
  \item review
  \item retrospective
  \item defined roles/accountabilities
\end{itemize}

\medskip

\subsection{What is version control?}

A system for tracking changes to source code and collaborating among developers.

Example:

\begin{quote}
Git\\
\end{quote}

\medskip

\subsection{Git vs GitHub?}

Git is a distributed version-control system.

GitHub is a platform for hosting and collaborating around Git repositories.

\medskip

\subsection{What is debugging?}

The process of identifying, understanding and fixing defects in software.

\medskip

\subsection{Verification vs validation?}

\textbf{Verification:}

\begin{quote}
Are we building the product correctly?\\
\end{quote}

\textbf{Validation:}

\begin{quote}
Are we building the correct product?\\
\end{quote}

\medskip

\subsection{Unit testing?}

Testing an individual function, method or small component in isolation.

\medskip

\subsection{Integration testing?}

Tests whether multiple components work correctly together.

\medskip

\subsection{System testing?}

Tests the complete integrated system against requirements.

\medskip

\subsection{Acceptance testing?}

Determines whether the system meets business/user acceptance criteria.

\medskip

\section{Computer Architecture}

\subsection{What is instruction cycle?}

Conceptually:

\begin{lstlisting}
Fetch
 $\downarrow$
Decode
 $\downarrow$
Execute
 $\downarrow$
Memory access if needed
 $\downarrow$
Write back
\end{lstlisting}

Actual implementations use pipelines and more complex microarchitectures.

\medskip

\subsection{What is pipelining?}

Pipelining overlaps different stages of multiple instructions.

Example:

\begin{lstlisting}
Instruction 1: Fetch $\rightarrow$ Decode $\rightarrow$ Execute
Instruction 2:         Fetch $\rightarrow$ Decode $\rightarrow$ Execute
Instruction 3:                 Fetch $\rightarrow$ Decode $\rightarrow$ Execute
\end{lstlisting}

This improves instruction throughput.

\medskip

\subsection{What is a pipeline hazard?}

A situation that prevents the next instruction from executing normally.

Types:

\begin{itemize}
  \item structural hazard
  \item data hazard
  \item control hazard
\end{itemize}

\medskip

\subsection{What is RISC?}

Reduced Instruction Set Computer.

It emphasizes a relatively simple instruction set and regular instruction execution.

\medskip

\subsection{What is CISC?}

Complex Instruction Set Computer.

It supports a richer and more varied instruction set.

x86 is historically associated with CISC-style ISA design.

\medskip

\subsection{What is clock speed?}

The frequency of the CPU clock, typically measured in GHz.

It is not by itself a complete measure of CPU performance.

\medskip

\subsection{Why doesn't higher GHz always mean faster CPU?}

Performance also depends on:

\begin{itemize}
  \item instructions per cycle
  \item architecture
  \item cache
  \item memory subsystem
  \item number of cores
  \item workload
  \item compiler optimization
\end{itemize}

\medskip

\subsection{What is multicore processing?}

A CPU contains multiple processing cores capable of executing multiple instruction streams concurrently.

\medskip

\section{Compiler \& Programming Languages}

\subsection{What is a compiler?}

A compiler translates source code into another form, commonly machine code or intermediate representation.

\medskip

\subsection{What is an interpreter?}

An interpreter executes a program through runtime interpretation rather than producing a conventional native executable as the primary execution artifact.

Modern implementations often combine interpretation with JIT compilation.

\medskip

\subsection{Compiler vs interpreter?}

Traditional compiler:

\begin{lstlisting}
Source $\rightarrow$ compiled program $\rightarrow$ execute
\end{lstlisting}

Traditional interpreter:

\begin{lstlisting}
Source $\rightarrow$ runtime execution
\end{lstlisting}

Modern language implementations frequently combine both techniques.

\medskip

\subsection{Compiler phases?}

A simplified compiler pipeline:

\begin{lstlisting}
Lexical Analysis
       $\downarrow$
Syntax Analysis
       $\downarrow$
Semantic Analysis
       $\downarrow$
Intermediate Representation
       $\downarrow$
Optimization
       $\downarrow$
Code Generation
\end{lstlisting}

\medskip

\subsection{What is lexical analysis?}

Converts character sequences into tokens.

Example:

\begin{lstlisting}
int x = 10;
\end{lstlisting}

Tokens include:

\begin{lstlisting}
int
x
=
10
;
\end{lstlisting}

\medskip

\subsection{What is parsing?}

Parsing analyzes token sequences according to the grammar of the language.

\medskip

\subsection{What is semantic analysis?}

Checks meaning-related rules.

Examples:

\begin{itemize}
  \item type compatibility
  \item variable declarations
  \item scope
  \item function arguments
\end{itemize}

\medskip

\subsection{Syntax error vs semantic error?}

Syntax error:

\begin{quote}
Violates language grammar.\\
\end{quote}

Semantic error:

\begin{quote}
Grammatically valid but violates language meaning/rules.\\
\end{quote}

\medskip

\section{Operating System + File Systems}

\subsection{What is a file system?}

A file system organizes and manages files and directories on storage devices.

Examples:

\begin{itemize}
  \item NTFS
  \item ext4
  \item APFS
  \item FAT32
\end{itemize}

\medskip

\subsection{What is a file?}

A named collection of data managed by the operating system/file system.

\medskip

\subsection{What is a directory?}

A structure used to organize references to files and other directories.

\medskip

\subsection{Absolute vs relative path?}

Absolute path starts from a root/location.

Relative path is interpreted relative to the current working directory.

\medskip

\subsection{What is a file descriptor?}

A file descriptor is a small integer handle used by Unix-like operating systems to refer to an open file or I/O resource.

Standard descriptors:

\begin{lstlisting}
0 $\rightarrow$ stdin
1 $\rightarrow$ stdout
2 $\rightarrow$ stderr
\end{lstlisting}

\medskip

\section{Cybersecurity}

\subsection{What is cybersecurity?}

Protection of systems, networks, applications and data against unauthorized access, disruption, modification or destruction.

\medskip

\subsection{CIA triad?}

\subsubsection{Confidentiality}

Only authorized parties can access information.

\subsubsection{Integrity}

Information is not improperly modified.

\subsubsection{Availability}

Authorized users can access resources when needed.

\medskip

\subsection{Authentication vs authorization?}

\textbf{Authentication:}

\begin{quote}
Who are you?\\
\end{quote}

\textbf{Authorization:}

\begin{quote}
What are you allowed to do?\\
\end{quote}

\medskip

\subsection{Encryption?}

Encryption transforms plaintext into ciphertext using a cryptographic process and key.

\medskip

\subsection{Symmetric encryption?}

The same secret key is used for encryption and decryption.

Example:

\begin{quote}
AES\\
\end{quote}

\medskip

\subsection{Asymmetric encryption?}

Uses a public/private key pair.

Examples:

\begin{itemize}
  \item RSA
  \item elliptic-curve cryptography
\end{itemize}

\medskip

\subsection{Hashing vs encryption?}

Hashing:

\begin{quote}
One-way transformation designed for integrity/fingerprinting.\\
\end{quote}

Encryption:

\begin{quote}
Reversible with the appropriate key.\\
\end{quote}

\medskip

\subsection{Why shouldn't passwords be stored using plain encryption?}

Password systems generally need to store password-derived verifiers rather than recoverable passwords.

Use a password-hashing scheme such as:

\begin{itemize}
  \item Argon2
  \item bcrypt
  \item scrypt
\end{itemize}

with appropriate salting and parameters.

\medskip

\subsection{What is SQL injection?}

An attack where untrusted input alters the intended SQL query structure.

Primary defense:

\begin{quote}
Parameterized queries/prepared statements.\\
\end{quote}

\medskip

\subsection{What is XSS?}

\textbf{Cross-Site Scripting} occurs when attacker-controlled content is executed as script in another user's browser.

Defenses include:

\begin{itemize}
  \item output encoding
  \item safe templating
  \item input validation where appropriate
  \item Content Security Policy
  \item avoiding unsafe HTML injection
\end{itemize}

\medskip

\subsection{What is CSRF?}

Cross-Site Request Forgery tricks an authenticated browser into making an unwanted state-changing request.

Common defenses:

\begin{itemize}
  \item CSRF tokens
  \item SameSite cookies
  \item appropriate origin checks
\end{itemize}

\medskip

\subsection{What is a firewall?}

A security control that filters network traffic according to defined rules.

\medskip

\subsection{What is malware?}

Malicious software.

Examples:

\begin{itemize}
  \item virus
  \item worm
  \item trojan
  \item ransomware
  \item spyware
\end{itemize}

\medskip

\subsection{Virus vs worm?}

A virus generally requires a host/program or user-assisted propagation mechanism.

A worm can self-propagate across networks.

\medskip

\section{Web \& Modern Application Architecture}

\subsection{Frontend vs backend?}

Frontend handles the user-facing interface.

Backend handles:

\begin{itemize}
  \item business logic
  \item data access
  \item authentication
  \item APIs
  \item server-side processing
\end{itemize}

\medskip

\subsection{What is a database server?}

A system running DBMS software that accepts database requests and manages stored data.

\medskip

\subsection{What is a web server?}

Software that handles HTTP requests and serves web resources or forwards requests to application logic.

Examples:

\begin{itemize}
  \item Apache
  \item Nginx
\end{itemize}

\medskip

\subsection{What is a reverse proxy?}

A server that receives requests on behalf of backend servers and forwards them appropriately.

It can provide:

\begin{itemize}
  \item TLS termination
  \item load balancing
  \item caching
  \item routing
  \item security controls
\end{itemize}

\medskip

\subsection{What is load balancing?}

Distributing incoming workloads across multiple servers.

Benefits:

\begin{itemize}
  \item scalability
  \item availability
  \item fault tolerance
\end{itemize}

\medskip

\subsection{Vertical vs horizontal scaling?}

\subsubsection{Vertical scaling}

Increase resources of one machine.

\begin{lstlisting}
More CPU/RAM
\end{lstlisting}

\subsubsection{Horizontal scaling}

Add more machines/instances.

\begin{lstlisting}
Server 1
Server 2
Server 3
\end{lstlisting}

\medskip

\section{Cloud Computing}

\subsection{What is cloud computing?}

On-demand delivery of computing resources over a network with elastic scaling and usage-based resource allocation.

\medskip

\subsection{IaaS?}

\textbf{Infrastructure as a Service}

Provides resources such as:

\begin{itemize}
  \item virtual machines
  \item storage
  \item networking
\end{itemize}

\medskip

\subsection{PaaS?}

\textbf{Platform as a Service}

Provides a managed platform for deploying applications.

\medskip

\subsection{SaaS?}

\textbf{Software as a Service}

Users consume a complete software application.

Examples include web-based productivity applications.

\medskip

\subsection{Public vs private cloud?}

Public cloud infrastructure is provided by a third-party cloud provider.

Private cloud is dedicated to a particular organization.

\medskip

\subsection{What is virtualization?}

Running virtualized computing environments on shared physical hardware through a hypervisor or similar mechanism.

\medskip

\subsection{What is a container?}

A lightweight isolated application environment that shares the host operating system kernel.

Docker is a common container technology.

\medskip

\subsection{Container vs virtual machine?}

VM:

\begin{lstlisting}
Hardware
 $\downarrow$
Hypervisor
 $\downarrow$
Guest OS
 $\downarrow$
Application
\end{lstlisting}

Container:

\begin{lstlisting}
Hardware
 $\downarrow$
Host OS
 $\downarrow$
Container runtime
 $\downarrow$
Containers
\end{lstlisting}

Containers are generally lighter because they share the host kernel.

\medskip

\section{Artificial Intelligence \& Machine Learning}

Because your background includes AI/ML, the board may probe this area.

\subsection{What is AI?}

Artificial Intelligence is the field of building systems capable of performing tasks that normally require aspects of human intelligence.

\medskip

\subsection{Machine learning?}

Machine learning enables systems to learn patterns from data rather than being explicitly programmed with every rule.

\medskip

\subsection{Deep learning?}

Deep learning uses multi-layer neural networks to learn hierarchical representations from data.

\medskip

\subsection{Supervised learning?}

Learning from labeled examples.

Examples:

\begin{itemize}
  \item classification
  \item regression
\end{itemize}

\medskip

\subsection{Unsupervised learning?}

Learning patterns from unlabeled data.

Examples:

\begin{itemize}
  \item clustering
  \item dimensionality reduction
\end{itemize}

\medskip

\subsection{Classification vs regression?}

Classification predicts discrete categories.

Regression predicts continuous numerical values.

\medskip

\subsection{What is overfitting?}

The model learns training-specific patterns, including noise, so training performance is high but generalization performance is poor.

\medskip

\subsection{How can overfitting be reduced?}

\begin{itemize}
  \item more data
  \item regularization
  \item dropout
  \item data augmentation
  \item early stopping
  \item simpler models
  \item cross-validation
\end{itemize}

\medskip

\subsection{What is underfitting?}

The model is too simple or insufficiently trained to capture the underlying pattern.

Both training and validation performance are poor.

\medskip

\subsection{Bias vs variance?}

\textbf{High bias:}

Model is too simple.

$\rightarrow$ underfitting.

\textbf{High variance:}

Model is overly sensitive to training data.

$\rightarrow$ overfitting.

\medskip

\subsection{What is a neural network?}

A computational model consisting of interconnected parameterized units that transform inputs through learned weights and nonlinear activation functions.

\medskip

\subsection{What is an activation function?}

A nonlinear function applied to a neuron's output.

Examples:

\begin{itemize}
  \item ReLU
  \item sigmoid
  \item tanh
  \item softmax
\end{itemize}

\medskip

\subsection{Why is ReLU popular?}

It is computationally simple and introduces nonlinearity while generally reducing the severe saturation problem of sigmoid in positive regions.

\medskip

\subsection{What is softmax?}

Softmax converts logits into a probability distribution over classes.

For class i:

\$\$
P\_i = \textbackslash\{\}frac\{e\textasciicircum{}\{z\_i\}\}\{\textbackslash\{\}sum\_j e\textasciicircum{}\{z\_j\}\}
\$\$

\medskip

\subsection{What is gradient descent?}

An optimization algorithm that updates parameters in the direction that reduces the loss.

Conceptually:

\$\$
\textbackslash\{\}theta \textbackslash\{\}leftarrow \textbackslash\{\}theta-\textbackslash\{\}eta\textbackslash\{\}nabla\_\textbackslash\{\}theta L
\$\$

where:

\begin{itemize}
  \item  = parameters
  \item  = learning rate
  \item L = loss
\end{itemize}

\medskip

\subsection{What is learning rate?}

The step size used by an optimization algorithm when updating model parameters.

Too large:

\begin{quote}
unstable/divergent training.\\
\end{quote}

Too small:

\begin{quote}
slow convergence.\\
\end{quote}

\medskip

\subsection{What is an epoch?}

One complete pass through the training dataset.

\medskip

\subsection{What is batch size?}

Number of training examples processed before a parameter update.

\medskip

\subsection{What is an iteration?}

One parameter-update step.

Approximately:

\$\$
\textbackslash\{\}text\{iterations per epoch\}
=
\textbackslash\{\}frac\{\textbackslash\{\}text\{training samples\}\}\{\textbackslash\{\}text\{batch size\}\}
\$\$

\medskip

\subsection{What is a confusion matrix?}

A table showing predicted vs actual class outcomes.

For binary classification:

\begin{lstlisting}
                 Predicted
              Positive Negative

Actual Pos       TP       FN
Actual Neg       FP       TN
\end{lstlisting}

\medskip

\subsection{Accuracy?}

\$\$
Accuracy=\textbackslash\{\}frac\{TP+TN\}\{TP+TN+FP+FN\}
\$\$

\medskip

\subsection{Precision?}

\$\$
Precision=\textbackslash\{\}frac\{TP\}\{TP+FP\}
\$\$

It answers:

\begin{quote}
Of predicted positives, how many were actually positive?\\
\end{quote}

\medskip

\subsection{Recall?}

\$\$
Recall=\textbackslash\{\}frac\{TP\}\{TP+FN\}
\$\$

It answers:

\begin{quote}
Of actual positives, how many did we identify?\\
\end{quote}

\medskip

\subsection{F1-score?}

Harmonic mean of precision and recall:

\$\$
F1=2\textbackslash\{\}frac\{Precision\textbackslash\{\}times Recall\}\{Precision+Recall\}
\$\$

\medskip

\subsection{Why can accuracy be misleading?}

Especially with imbalanced datasets.

Example:

\begin{lstlisting}
99% healthy
1% diseased
\end{lstlisting}

A model predicting "healthy" every time gets 99\% accuracy but detects no disease.

\medskip

\subsection{What is transfer learning?}

Using knowledge learned by a model on one task/dataset as a starting point for another related task.

\medskip

\subsection{What is fine-tuning?}

Continuing training of a pretrained model, usually with a smaller learning rate, so that its parameters adapt to the target task.

\medskip

\section{Software Design}

\subsection{What is modularity?}

Dividing software into relatively independent modules with clear responsibilities and interfaces.

\medskip

\subsection{What is cohesion?}

How strongly related the responsibilities inside a module are.

High cohesion is generally desirable.

\medskip

\subsection{What is coupling?}

How strongly modules depend on one another.

Low coupling is generally desirable.

\medskip

\subsection{High cohesion and low coupling --- why desirable?}

It makes software:

\begin{itemize}
  \item easier to understand
  \item easier to test
  \item easier to modify
  \item easier to maintain
\end{itemize}

\medskip

\subsection{What is SOLID?}

Five object-oriented design principles:

\begin{lstlisting}
S $\rightarrow$ Single Responsibility
O $\rightarrow$ Open/Closed
L $\rightarrow$ Liskov Substitution
I $\rightarrow$ Interface Segregation
D $\rightarrow$ Dependency Inversion
\end{lstlisting}

\medskip

\subsection{Single Responsibility Principle?}

A class/module should have one coherent responsibility and one primary reason to change.

\medskip

\subsection{Open/Closed Principle?}

Software entities should be:

\begin{quote}
open for extension, closed for modification.\\
\end{quote}

\medskip

\subsection{Liskov Substitution Principle?}

Subtypes should be substitutable for their base types without breaking expected behavior.

\medskip

\subsection{Dependency Inversion Principle?}

High-level modules should not depend directly on low-level implementation details.

Both should depend on abstractions.

\medskip

\section{Git \& Development}

\subsection{What is Git?}

A distributed version-control system.

\medskip

\subsection{What is a commit?}

A recorded snapshot of changes in a Git repository.

\medskip

\subsection{What is a branch?}

An independent line of development.

\medskip

\subsection{What is merge?}

Combining changes from different branches.

\medskip

\subsection{What is a merge conflict?}

When Git cannot automatically reconcile conflicting changes.

\medskip

\subsection{What is git clone?}

Copies an existing repository to your local system.

\medskip

\subsection{What is git pull?}

Fetches remote changes and integrates them into the current branch, depending on configuration.

\medskip

\subsection{What is git fetch?}

Downloads changes from a remote repository without automatically integrating them into the current branch.

\medskip

\section{Common Programming Viva Traps}

\subsection{What is the difference between = and ==?}

\texttt{=} is assignment.

\texttt{==} compares equality.

\medskip

\subsection{What is ++i vs i++?}

Both increment \texttt{i}.

But:

\begin{lstlisting}
++i
\end{lstlisting}

increments before yielding the value.

\begin{lstlisting}
i++
\end{lstlisting}

yields the old value before incrementing.

\medskip

\subsection{What is pass-by-value?}

A copy of the argument is passed to the function.

Changes to the parameter don't directly modify the original variable.

\medskip

\subsection{What is pass-by-reference?}

The function receives an alias/reference to the original object, allowing modification of the original depending on language semantics.

\medskip

\subsection{What is scope?}

Scope determines where a name can be accessed in a program.

\medskip

\subsection{What is lifetime?}

Lifetime is the period during which an object or variable exists.

Scope and lifetime are related but not identical.

\medskip

\subsection{What is static variable?}

A static variable has a lifetime extending across the relevant program execution scope, while its accessibility depends on where it is declared.

In C/C++, \texttt{static} also affects linkage for certain file-scope declarations.

\medskip

\subsection{What is constant?}

A value that the program should not modify after initialization, depending on the language's const/final semantics.

\medskip

\section{Important Conceptual Comparisons}

These are \textbf{high-value viva questions} because interviewers love comparisons.

\subsection{Compiler vs interpreter}

\begin{lstlisting}
Compiler:
Source $\rightarrow$ translation $\rightarrow$ executable/IR

Interpreter:
Program $\rightarrow$ runtime execution
\end{lstlisting}

Modern systems often combine compilation and interpretation.

\medskip

\subsection{Stack vs heap}

\begin{lstlisting}
Stack:
Automatic/local lifetime
Fast
Structured by function calls

Heap:
Dynamic allocation
Flexible lifetime
Requires management/GC
\end{lstlisting}

\medskip

\subsection{TCP vs UDP}

\begin{lstlisting}
TCP:
Reliable + ordered + connection-oriented

UDP:
Minimal + connectionless + best-effort
\end{lstlisting}

\medskip

\subsection{Process vs thread}

\begin{lstlisting}
Process:
Independent address space

Thread:
Execution unit within process
\end{lstlisting}

\medskip

\subsection{Primary key vs foreign key}

\begin{lstlisting}
Primary key:
Identifies row

Foreign key:
References another table's key
\end{lstlisting}

\medskip

\subsection{WHERE vs HAVING}

\begin{lstlisting}
WHERE:
Filters rows

HAVING:
Filters groups
\end{lstlisting}

\medskip

\subsection{DELETE vs TRUNCATE vs DROP}

\begin{lstlisting}
DELETE:
Rows

TRUNCATE:
All rows

DROP:
Entire database object
\end{lstlisting}

\medskip

\subsection{Array vs linked list}

\begin{lstlisting}
Array:
Fast indexing

Linked list:
Flexible node-based insertion/deletion
\end{lstlisting}

\medskip

\subsection{BFS vs DFS}

\subsubsection{BFS}

Uses:

\begin{quote}
Queue\\
\end{quote}

Useful for shortest paths in \textbf{unweighted graphs}.

\subsubsection{DFS}

Uses:

\begin{quote}
Stack/recursion\\
\end{quote}

Useful for:

\begin{itemize}
  \item traversal
  \item cycle detection
  \item topological algorithms
  \item connected components
\end{itemize}

\medskip

\subsection{What is BFS complexity?}

For adjacency-list representation:

\$\$
O(V+E)
\$\$

\medskip

\subsection{DFS complexity?}

Also:

\$\$
O(V+E)
\$\$

with adjacency lists.

\medskip

\subsection{What is a tree?}

A connected acyclic graph.

For a tree with n vertices:

\$\$
E=n-1
\$\$

\medskip

\subsection{What is a binary tree?}

A tree where each node has at most two children.

\medskip

\subsection{What is a binary search tree?}

A binary tree satisfying an ordering property such as:

\begin{lstlisting}
left subtree < node < right subtree
\end{lstlisting}

for unique keys in the standard formulation.

\medskip

\subsection{BST search complexity?}

Average:

\begin{lstlisting}
O(log n)
\end{lstlisting}

for a reasonably balanced tree.

Worst case:

\begin{lstlisting}
O(n)
\end{lstlisting}

if the tree becomes highly skewed.

\medskip

\subsection{What is AVL tree?}

A self-balancing binary search tree where the balance factor of each node is restricted, commonly to -1, 0 or +1.

Search/insert/delete:

\begin{lstlisting}
O(log n)
\end{lstlisting}

\medskip

\subsection{What is a heap?}

A complete binary tree satisfying a heap-order property.

Max heap:

\begin{lstlisting}
parent $\geq$ children
\end{lstlisting}

Min heap:

\begin{lstlisting}
parent $\leq$ children
\end{lstlisting}

\medskip

\subsection{Heap applications?}

\begin{itemize}
  \item priority queues
  \item heap sort
  \item scheduling
  \item graph algorithms such as Dijkstra/Prim implementations
\end{itemize}

\medskip

\section{Advanced CS Questions}

\subsection{What is dynamic programming?}

Dynamic programming solves problems with overlapping subproblems and optimal substructure by storing previously computed results.

Two common approaches:

\begin{itemize}
  \item memoization
  \item tabulation
\end{itemize}

\medskip

\subsection{Greedy vs dynamic programming?}

Greedy makes the locally best choice at each step.

Dynamic programming systematically considers subproblem solutions.

Greedy works only when the problem has the appropriate greedy-choice property.

\medskip

\subsection{What is divide and conquer?}

Break a problem into smaller independent subproblems, solve them recursively, and combine their solutions.

Example:

\begin{quote}
Merge sort.\\
\end{quote}

\medskip

\subsection{What is a graph?}

A graph consists of:

\begin{itemize}
  \item vertices/nodes
  \item edges
\end{itemize}

It can model relationships such as:

\begin{lstlisting}
Cities $\leftrightarrow$ Roads
Users $\leftrightarrow$ Friendships
Computers $\leftrightarrow$ Network links
\end{lstlisting}

\medskip

\subsection{Directed vs undirected graph?}

Directed:

\begin{lstlisting}
A $\rightarrow$ B
\end{lstlisting}

Undirected:

\begin{lstlisting}
A --- B
\end{lstlisting}

\medskip

\subsection{Weighted graph?}

Edges have associated numerical weights.

Example:

\begin{lstlisting}
Dhaka --5km-- Gazipur
\end{lstlisting}

\medskip

\subsection{Dijkstra's algorithm?}

Finds shortest paths from a source in graphs with \textbf{non-negative edge weights}.

\medskip

\subsection{Can Dijkstra handle negative edges?}

No.

Bellman-Ford can handle negative edge weights and can detect reachable negative cycles.

\medskip

\subsection{What is a minimum spanning tree?}

A spanning tree of a connected weighted undirected graph with minimum total edge weight.

Algorithms:

\begin{itemize}
  \item Kruskal
  \item Prim
\end{itemize}

\medskip

\subsection{What is topological sorting?}

A linear ordering of vertices in a directed acyclic graph such that:

\begin{lstlisting}
u $\rightarrow$ v
\end{lstlisting}

means u appears before v.

\medskip

\subsection{Where is topological sorting used?}

\begin{itemize}
  \item course prerequisites
  \item build systems
  \item dependency management
  \item task scheduling
\end{itemize}

\medskip

\section{Distributed Systems}

\subsection{What is a distributed system?}

A collection of independent computers that cooperate and appear to users as a coherent system.

\medskip

\subsection{Why use distributed systems?}

\begin{itemize}
  \item scalability
  \item fault tolerance
  \item geographic distribution
  \item availability
  \item performance
\end{itemize}

\medskip

\subsection{What is replication?}

Maintaining multiple copies of data/services.

Benefits:

\begin{itemize}
  \item availability
  \item fault tolerance
  \item read scalability
\end{itemize}

\medskip

\subsection{What is sharding?}

Partitioning data across multiple machines/nodes.

\medskip

\subsection{Replication vs sharding?}

Replication:

\begin{quote}
Same data copied to multiple nodes.\\
\end{quote}

Sharding:

\begin{quote}
Different portions of data distributed among nodes.\\
\end{quote}

They can be combined.

\medskip

\subsection{What is CAP theorem?}

In the presence of a network partition, a distributed system cannot simultaneously guarantee both:

\begin{itemize}
  \item strong consistency
  \item full availability
\end{itemize}

while also tolerating the partition.

The three properties are:

\begin{lstlisting}
Consistency
Availability
Partition tolerance
\end{lstlisting}

\medskip

\subsection{What is eventual consistency?}

If no new updates occur and communication continues, replicas will eventually converge to the same value.

\medskip

\section{Modern Technology Questions}

\subsection{What is cloud computing?}

On-demand computing resources delivered over networks with elasticity and resource pooling.

\medskip

\subsection{What is edge computing?}

Processing data closer to where it is generated instead of sending everything to a centralized cloud.

Benefits:

\begin{itemize}
  \item lower latency
  \item reduced bandwidth
  \item improved local responsiveness
\end{itemize}

\medskip

\subsection{What is IoT?}

Internet of Things refers to interconnected physical devices capable of sensing, processing and communicating data.

Examples:

\begin{itemize}
  \item smart meters
  \item industrial sensors
  \item wearable devices
\end{itemize}

\medskip

\subsection{What is blockchain?}

A distributed ledger technology in which transactions are recorded using cryptographic mechanisms and consensus protocols.

\medskip

\subsection{What is generative AI?}

AI systems capable of generating new content such as:

\begin{itemize}
  \item text
  \item images
  \item audio
  \item video
  \item code
\end{itemize}

\medskip

\subsection{What is an LLM?}

Large Language Model.

A neural language model trained on large datasets to model and generate natural language.

Modern LLMs are commonly based on Transformer architectures.

\medskip

\subsection{What is a Transformer?}

A neural-network architecture based heavily on attention mechanisms rather than recurrence.

It powers many modern language and multimodal models.

\medskip

\subsection{What is attention?}

Attention allows a model to assign different weights to different parts of the input when producing a representation or output.

\medskip

\subsection{What is Generative AI vs traditional ML?}

Traditional ML often predicts or classifies:

\begin{lstlisting}
Input $\rightarrow$ Prediction
\end{lstlisting}

Generative AI can produce new content:

\begin{lstlisting}
Prompt/Input $\rightarrow$ Generated content
\end{lstlisting}

\medskip

\section{High-Probability Technical Questions}

These are the questions I would give \textbf{extra priority} because of the position.

\subsection{Why is data structure important for a programmer?}

Because algorithm performance depends heavily on how data is organized.

Choosing the right structure can turn an operation from:

\begin{lstlisting}
O(n)
\end{lstlisting}

into:

\begin{lstlisting}
O(1)
\end{lstlisting}

average-case, for example with hash-based lookup.

\medskip

\subsection{Which language would you use for system-level programming?}

C/C++ are common choices because they provide:

\begin{itemize}
  \item low-level control
  \item predictable performance
  \item direct memory manipulation
\end{itemize}

\medskip

\subsection{Which language would you choose for enterprise application development?}

It depends on requirements.

Common choices include:

\begin{itemize}
  \item Java
  \item C\#
  \item C++
  \item JavaScript/TypeScript
  \item Python
\end{itemize}

\medskip

\subsection{How would you design an employee management system?}

I would first identify:

\subsubsection{Entities}

\begin{lstlisting}
Employee
Department
Position
Salary
Attendance
\end{lstlisting}

Then define:

\begin{itemize}
  \item relationships
  \item primary/foreign keys
  \item business rules
  \item API requirements
  \item authentication/authorization
  \item indexes
  \item transaction boundaries
\end{itemize}

\medskip

\subsection{How would you secure an organization's employee database?}

I would use:

\begin{itemize}
  \item authentication
  \item role-based authorization
  \item parameterized SQL
  \item encryption in transit
  \item encryption at rest where appropriate
  \item password hashing
  \item audit logs
  \item backups
  \item least privilege
  \item input validation
  \item regular security updates
\end{itemize}

\medskip

\subsection{How would you improve a slow database query?}

I would investigate:

\begin{enumerate}
  \item Execution plan
  \item Missing/ineffective indexes
  \item Join conditions
  \item Filtering
  \item Excessive data retrieval
  \item Sorting/grouping
  \item Statistics
  \item Query structure
\end{enumerate}

Then benchmark changes rather than assuming an index will always help.

\medskip

\subsection{How would you troubleshoot a slow application?}

I'd follow a measurement-driven process:

\begin{lstlisting}
Reproduce
   $\downarrow$
Measure
   $\downarrow$
Identify bottleneck
   $\downarrow$
Optimize
   $\downarrow$
Test
   $\downarrow$
Monitor
\end{lstlisting}

Potential bottlenecks:

\begin{itemize}
  \item CPU
  \item memory
  \item disk I/O
  \item database
  \item network
  \item inefficient algorithms
  \item external APIs
\end{itemize}

\medskip

\subsection{What would you do if a program suddenly crashes?}

I'd:

\begin{enumerate}
  \item Reproduce the issue.
  \item Inspect logs/error messages.
  \item Identify the failing component.
  \item Use a debugger if appropriate.
  \item Check memory/resource issues.
  \item Fix the root cause.
  \item Add regression tests.
  \item Verify the fix.
\end{enumerate}

\medskip

\subsection{What is the most important skill of a programmer?}

I would answer:

\begin{quote}
\textbf{Problem-solving ability.}\\
\end{quote}

Programming languages and frameworks can change, but the ability to decompose a problem, design an algorithm, reason about complexity, debug and verify a solution remains fundamental.

\medskip

\section{Very Short Rapid-Fire Round}

These are designed for the situation where the board fires questions one after another.

\subsubsection{What is an algorithm?}

A finite sequence of steps for solving a problem.

\subsubsection{What is a data structure?}

A method of organizing data for efficient operations.

\subsubsection{Stack principle?}

LIFO.

\subsubsection{Queue principle?}

FIFO.

\subsubsection{Binary search complexity?}

O(log n).

\subsubsection{Linear search complexity?}

O(n).

\subsubsection{Merge sort?}

O(n log n).

\subsubsection{Worst-case quicksort?}

O(n\textbackslash\{\}textsuperscript\{2\}).

\subsubsection{Hash table average lookup?}

O(1).

\subsubsection{Tree with no cycles?}

Tree.

\subsubsection{Tree with two children maximum?}

Binary tree.

\subsubsection{Sorted binary tree?}

BST.

\subsubsection{Graph shortest path with nonnegative weights?}

Dijkstra.

\subsubsection{Graph traversal using queue?}

BFS.

\subsubsection{Graph traversal using stack?}

DFS.

\subsubsection{SQL full form?}

Structured Query Language.

\subsubsection{DBMS full form?}

Database Management System.

\subsubsection{Primary key?}

Unique row identifier.

\subsubsection{Foreign key?}

Reference to another table's key.

\subsubsection{ACID?}

Atomicity, Consistency, Isolation, Durability.

\subsubsection{OS full form?}

Operating System.

\subsubsection{CPU?}

Central Processing Unit.

\subsubsection{RAM?}

Random Access Memory.

\subsubsection{ROM?}

Read Only Memory.

\subsubsection{HTTP?}

HyperText Transfer Protocol.

\subsubsection{HTTPS?}

HTTP over TLS.

\subsubsection{IP?}

Internet Protocol.

\subsubsection{TCP?}

Transmission Control Protocol.

\subsubsection{UDP?}

User Datagram Protocol.

\subsubsection{DNS?}

Domain Name System.

\subsubsection{DHCP?}

Dynamic Host Configuration Protocol.

\subsubsection{LAN?}

Local Area Network.

\subsubsection{WAN?}

Wide Area Network.

\subsubsection{API?}

Application Programming Interface.

\subsubsection{OOP pillars?}

Encapsulation, abstraction, inheritance, polymorphism.

\subsubsection{Git?}

Distributed version-control system.

\subsubsection{Compiler?}

Translates source code into another executable/intermediate representation.

\subsubsection{Deadlock?}

Processes wait indefinitely for resources held by one another.

\subsubsection{Race condition?}

Outcome depends on concurrent execution timing.

\subsubsection{Semaphore?}

Synchronization/resource-counting primitive.

\subsubsection{Mutex?}

Mutual-exclusion lock.

\subsubsection{Virtual memory?}

Abstraction providing processes with virtual address spaces backed by physical memory/storage.

\subsubsection{Page fault?}

Access to a page not currently resident in physical memory.

\subsubsection{Cache?}

Small, fast memory storing frequently/recently used data.

\subsubsection{AI?}

Artificial Intelligence.

\subsubsection{ML?}

Machine Learning.

\subsubsection{CNN?}

Convolutional Neural Network.

\subsubsection{RNN?}

Recurrent Neural Network.

\subsubsection{Transformer?}

Attention-based neural architecture.

\subsubsection{Overfitting?}

Excellent training fit but poor generalization.

\medskip

\section{Questions the Board May Ask After Your Answer}

This is particularly important.

Suppose they ask:

\begin{quote}
\textbf{``What is a primary key?''}\\
\end{quote}

Don't assume the conversation ends there.

They may immediately ask:

\begin{quote}
Can it be NULL?\\
\end{quote}

\textbf{No.}

\begin{quote}
Can a table have two primary keys?\\
\end{quote}

\textbf{No. A table has one primary-key constraint, although it may contain multiple columns as a composite primary key.}

Then:

\begin{quote}
What is a candidate key?\\
\end{quote}

\textbf{A minimal attribute set that uniquely identifies a row.}

Then:

\begin{quote}
What is a super key?\\
\end{quote}

\textbf{Any attribute set that uniquely identifies a row; it need not be minimal.}

Then:

\begin{quote}
What is a foreign key?\\
\end{quote}

\textbf{An attribute or set of attributes referencing a key in another or sometimes the same table.}

That is how a viva can escalate.

\medskip

\section{Follow-up Chains You Should Practice}

\subsection{Chain 1: Array}

\begin{quote}
What is an array?\\
\end{quote}

$\rightarrow$ Contiguous collection.

\begin{quote}
Why is access O(1)?\\
\end{quote}

$\rightarrow$ Address calculation.

\begin{quote}
Why is insertion O(n)?\\
\end{quote}

$\rightarrow$ Elements may need shifting.

\begin{quote}
What if insertion is at the end?\\
\end{quote}

$\rightarrow$ O(1) amortized for dynamic arrays.

\begin{quote}
Why is vector faster than linked list in many workloads?\\
\end{quote}

$\rightarrow$ Better cache locality.

\medskip

\section{Chain 2: Database}

\begin{quote}
What is normalization?\\
\end{quote}

$\rightarrow$ Reducing redundancy/dependency anomalies.

\begin{quote}
Why normalize?\\
\end{quote}

$\rightarrow$ Avoid insertion/update/deletion anomalies.

\begin{quote}
What is 1NF?\\
\end{quote}

$\rightarrow$ Atomic attributes/no repeating groups.

\begin{quote}
2NF?\\
\end{quote}

$\rightarrow$ No partial dependency.

\begin{quote}
3NF?\\
\end{quote}

$\rightarrow$ No relevant transitive dependency of non-key attributes.

\begin{quote}
Can normalization hurt performance?\\
\end{quote}

$\rightarrow$ Yes; excessive joins may increase read cost.

\medskip

\section{Chain 3: Networking}

\begin{quote}
What is TCP?\\
\end{quote}

$\rightarrow$ Reliable connection-oriented transport protocol.

\begin{quote}
How does TCP establish a connection?\\
\end{quote}

$\rightarrow$ Three-way handshake.

\begin{lstlisting}
Client $\rightarrow$ SYN
Server $\rightarrow$ SYN-ACK
Client $\rightarrow$ ACK
\end{lstlisting}

\begin{quote}
Why three messages?\\
\end{quote}

$\rightarrow$ To establish and confirm both sides' sequence/acknowledgment state.

\begin{quote}
TCP vs UDP?\\
\end{quote}

$\rightarrow$ Reliability/ordering/control versus lightweight best-effort transport.

\medskip

\section{Chain 4: Operating System}

\begin{quote}
What is a process?\\
\end{quote}

$\rightarrow$ Program in execution.

\begin{quote}
Process vs thread?\\
\end{quote}

$\rightarrow$ Separate address space vs shared process resources.

\begin{quote}
What is context switching?\\
\end{quote}

$\rightarrow$ Save one execution state and restore another.

\begin{quote}
What is deadlock?\\
\end{quote}

$\rightarrow$ Circular resource waiting.

\begin{quote}
Four conditions?\\
\end{quote}

$\rightarrow$ Mutual exclusion, hold-and-wait, no preemption, circular wait.

\medskip

\section{Chain 5: OOP}

\begin{quote}
What is OOP?\\
\end{quote}

$\rightarrow$ Object-centered programming paradigm.

\begin{quote}
Four pillars?\\
\end{quote}

$\rightarrow$ Encapsulation, abstraction, inheritance, polymorphism.

\begin{quote}
Encapsulation?\\
\end{quote}

$\rightarrow$ Controlled access to state and behavior bundling.

\begin{quote}
Abstraction?\\
\end{quote}

$\rightarrow$ Hide unnecessary implementation details.

\begin{quote}
Overloading?\\
\end{quote}

$\rightarrow$ Same method name, different parameter lists.

\begin{quote}
Overriding?\\
\end{quote}

$\rightarrow$ Subclass replaces inherited implementation.

\medskip

\section{Chain 6: Algorithm}

\begin{quote}
What is binary search?\\
\end{quote}

$\rightarrow$ Search sorted data by repeatedly halving the search space.

\begin{quote}
Complexity?\\
\end{quote}

$\rightarrow$ O(log n).

\begin{quote}
Why?\\
\end{quote}

$\rightarrow$ Search space halves each step.

\begin{quote}
Can it work on unsorted data?\\
\end{quote}

$\rightarrow$ Not in its standard form.

\begin{quote}
What is interpolation search?\\
\end{quote}

$\rightarrow$ Uses estimated position based on key distribution; can outperform binary search on uniformly distributed data but has worse worst-case behavior.

\medskip

\section{Chain 7: Security}

\begin{quote}
What is authentication?\\
\end{quote}

$\rightarrow$ Verifying identity.

\begin{quote}
Authorization?\\
\end{quote}

$\rightarrow$ Determining permissions.

\begin{quote}
SQL injection?\\
\end{quote}

$\rightarrow$ Input manipulates SQL query structure.

\begin{quote}
Defense?\\
\end{quote}

$\rightarrow$ Prepared statements/parameterized queries.

\begin{quote}
Password storage?\\
\end{quote}

$\rightarrow$ Salted password hashing using an appropriate password-hashing algorithm.

\medskip